% Generated draft: requires mathematical and source review.
\documentclass[11pt]{article}
\usepackage[margin=25mm]{geometry}
\usepackage{fontspec}
\setmainfont{Latin Modern Roman}
\usepackage{amsmath,amssymb,mathtools,unicode-math}
\setmathfont{Latin Modern Math}
\usepackage{enumitem,needspace,xurl,hyperref}
\setlength{\parindent}{0pt}
\setlength{\parskip}{5pt}
\setlength{\emergencystretch}{3em}
\newcommand{\N}{\mathbb{N}}
\newcommand{\Z}{\mathbb{Z}}
\newcommand{\Q}{\mathbb{Q}}
\newcommand{\R}{\mathbb{R}}
\newcommand{\C}{\mathbb{C}}
\newcommand{\F}{\mathbb{F}}
\newcommand{\A}{\mathbb{A}}
\newcommand{\PP}{\mathbb{P}}
\newcommand{\E}{\mathbb{E}}
\newcommand{\eps}{\varepsilon}
\newcommand{\dd}{\,\mathrm{d}}
\begin{document}
\section*{UnsolvedMath ID: 30005537}
\textbf{Source problem code:} catalogue entry
\subsection{Definitions and mathematical statement}
Let \(\mathbb H^n\) denote the Heisenberg group of real dimension \(2n+1\), with its standard CR structure, and write its homogeneous dimension as \(Q=2n+2\). The usual critical-exponent notation consistent with the supplied, typographically corrupted excerpt is
\[
q_* = \frac{Q+2}{Q-2},\qquad q_*+1=\frac{2Q}{Q-2}.
\]
Here \(n\) is a positive integer; the dimensions covered by the cited theorem and by each requested extension are not specified in the supplied record.

\textbf{Classification question, with unresolved source parameters.} Classify all positive solutions on the whole group \(\mathbb H^n\) of the critical CR Yamabe equation referred to as equation (1) in the source, without imposing the finite-volume hypothesis appearing in the cited theorem of D. Jerison and J. M. Lee, source reference [6]. More formally, let \(\mathcal A_n\) denote the source's admissible solution class and let \(\mathcal E_n(u)=0\) denote that critical equation. The requested object is a complete description of
\[
\mathcal S_n=\left\{u\in\mathcal A_n:\ u>0\text{ on }\mathbb H^n,\ \mathcal E_n(u)=0\text{ on }\mathbb H^n\right\},
\]
with no finite-volume restriction. The notation \(\mathcal A_n\) and \(\mathcal E_n\) records missing source information; it does not supply a definition of the absent equation or its solution concept. The original question asks whether the finite-volume hypothesis can be removed from the Jerison--Lee theorem, but neither that theorem's full hypotheses and conclusion nor its displayed solution family are included in the supplied excerpt. Consequently, an explicit quantified extension of that theorem cannot be recovered from this record alone.

The source additionally requests analogous classifications of positive solutions for fourth-order equations related to CR \(Q\)-curvature and for a \(p\)-Laplace equation on \(\mathbb H^n\). These are separate unresolved formulation targets: the fourth-order operator, equations, dimensional restrictions, \(p\)-range, right-hand sides, solution concepts, and any admissibility or asymptotic conditions are not supplied. No particular equations or additional boundary conditions are selected here.

\textbf{Supplied dated qualification.} The record's literature assessment, dated 2026-08-22, states that the positive-solution classification for the critical CR Yamabe equation without finite volume has been obtained, while the fourth-order CR \(Q\)-curvature and \(p\)-Laplace questions remain open. This is supplied dated curation, not an independently verified mathematical or literature assessment.
\subsection{Short English statement}
The original question asks whether the finite-volume assumption in the Jerison--Lee classification of positive critical CR Yamabe solutions on \(\mathbb H^n\) can be removed, and asks for analogous fourth-order CR \(Q\)-curvature and \(p\)-Laplace classifications. The supplied assessment dated 2026-08-22 labels the first part solved and the extensions open. The exact equations and several essential hypotheses are missing, so this statement retains those gaps explicitly and requires human mathematical review.
\subsection{Sources}
\begin{itemize}
\item Local UnsolvedMath registry, supplied complete record: permanent ID \(30005537\), title \emph{Classification of Critical CR Yamabe Solutions}, exact source problem code \texttt{catalogue entry}. The record attributes the problem to the 2023 Oberwolfach report; no individual proposer is supplied.
\item \emph{Hypoelliptic Operators in Geometry} (2023), Oberwolfach Reports, pp.\ 1323--1396, DOI \texttt{10.4171/owr/2023/24}: \url{https://doi.org/10.4171/owr/2023/24}. This is the authoritative report citation supplied by the registry; it has not been independently consulted for this draft.
\item Supplied literature source labeled ``OWR 2023 24'': \url{https://ems.press/content/serial-article-files/47017}.
\item Supplied literature source labeled ``Liu'': \url{https://arxiv.org/abs/2512.22458}. No paper title, full author name, theorem statement, or detailed scope is supplied.
\item The original excerpt attributes the finite-volume theorem to D. Jerison and J. M. Lee as source reference [6], and associates the CR Sobolev embedding exponent with source reference [2]. The corresponding bibliography entries are absent from the supplied record.
\item Literature checked at: \texttt{2026-08-22}. The supplied dated curation reports a classification without the finite-volume assumption for the critical CR Yamabe equation and describes the fourth-order CR \(Q\)-curvature and \(p\)-Laplace questions as open. These claims and the date are preserved as supplied curation, not independently verified findings. The registry labels the overall problem \texttt{partially\_solved} and the statement \texttt{corrected\_verified}, with a supplied assertion of reconstruction and checking against the official report; this draft does not independently establish those labels.
\end{itemize}
\end{document}
